\section{Introduction}
\label{section:intro}

In a multilinguistic society with a historically dominant language serving as the lingua franca, it is unclear whether the introduction of mother tongue instruction (MTI) increases employment opportunities.  Links between MTI and success in the labor market are theoretically unclear, nor empirically well-established, although whether MTI contributes to employment outcomes is a topical research question.  

Generally, MTI affects labor market outcomes through its effects on education and proficiency of a prominent language. Several studies document that MTI in early grades improves educational outcomes \citep[e.g.,][]{Walter2012,taylor_estimating_2016}. However, MTI does not necessarily lead to improved education, as the mediating factors can enhance or diminish access to curriculum content and learning \citep{Heugh2007}. The capacity of social groups to make use of their rights, language teaching methodologies, language teaching programs design, and availability of textbooks are some of the factors that determine the efficacy of learning in any language, suggesting that the impact of MTI on education may vary according to the context in which it is implemented. 

MTI may improve education and still result in poor job outcomes if it reduces proficiency in a dominant language. Although it helps to master international languages better at higher levels \citep[e.g.,][]{seid_impact_2019}, MTI can negatively affect employment and earnings by hampering fluency in a dominant language \citep[e.g.,][]{chiswick_endogeneity_2015}. Hence, whether MTI affects employment or not is an empirical research question with important policy implications. 

We test a conjecture that MTI induces differences in employment outcomes through both the education and the prominent-language channels, by exploiting a major policy change as a source of variation in the treatment. After installing ethnic federalism in 1994, Ethiopia introduced MTI in its primary schools, with each state adopting the national policy with different rates of intensity.\footnote{A map of Ethiopia divided into ethnic states is shown in Figure\ref{fig:map}.} Prior to 1994, Amharic was the language of instruction in elementary schools throughout the country, serving both as a source of contestation (for cultural reasons) and as a ladder of upward mobility in the center. Thus, a study based on the institution of MTI in this context may reveal interesting new insights on the relationships between the choice of a language of instruction and the outcomes of employment.  


We instrument variations in MTI intake by the interaction of ethnic identity of students and the MTI policies of the regional states, generating an exogenous IV named Predicted MTI. A review of the MTI adoptation processes in each state reveals that the geographic spread of MTI was randomly determined without the appropriate justifications \citep{Heugh2007}. In states where MTI is used less intensely than was recommended by national policy, for example, policy makers relied on unsubstantiated claims that English would be more beneficial as a medium of instruction in elementary schools \citep{Heugh2007}, although the dissemination and use of English is distinctly and similarly low in all federated states of Ethiopia \citep{gerencheal2019foreign}. 

We conduct two tests that suggest that Predicted MTI is as good as randomly assigned and satisfies the exclusion restriction. First, we show that there are no discernible correlations between employment and predicted MTI in samples where there is no obvious reason for associations between the treatment and employment outcomes. We also compare employment outcomes in states that implemented the MTI policy similarly to check if the treatment effects are driven by unobserved differences in state characteristics, which may lead to variations in demand for formal education and labor market outcomes.

We find that the impact of MTI on employment is mixed (Table \ref{tab:mainreg}).  MTI increased salary employment, but has no discernible effect on wage employment. The increase in salary employment in MTI is consistent with the evidence in Table \ref{tab:rev} that cognitive ability increased with MTI, suggesting that mastery of the dominant language, Amharic, is less important to obtain salaried jobs. 

The finding that the introduction of MTI in Ethiopia has proven to be a double-edged sword in terms of employment generation is reasonable considering the prevailing social infrastructure in the country. Most salaried jobs are in the public sector, where employees have legal rights to work in their native languages. In the private sector, where the mastery of Amharic is still essential to obtain wage employment, the positive effects of the treatment on human capital have been offset by its likely unintended adverse effects on the mastery of the dominant language (Table \ref{tab:rev}). Therefore, MTI increases employment outcomes only if it is complemented with appropriate institutional and policy measures.

These findings do not change when we exclude the Amhara state from the analysis. The magnitudes of MTI's effects on employment and the qualitative conclusions of the study are remarkably similar when the analysis is conducted using the overall sample or only the sample from states with previously suppressed languages, indicating that being from the region with a historically dominant language as such does not confer special advantages in the labor market. 

Our paper contributes to the emerging empirical literature that examines the effects of MTI on labor market outcomes. Studies that have found that MTI improves labor market outcomes link these gains to the positive effects of MTI on education \citep[e.g.,][]{eriksson_does_2014,seid_does_2016,Alidou2006,Walter2012,eriksson_does_2014,taylor_estimating_2016,seid_does_2016,piper_implementing_2016,ramachandran_language_2017,laitin_legacy_2019}, typically through increased participation in schools, reduced grade repetition and dropouts \citep{Benson2000,Benson2005,bender_their_2005} and increased chances of building non-language cognitive skills such as literacy and numeracy \citep{eriksson_does_2014,Trudell2012}.  

%The effectiveness of MTI is highly sensitive to choices of inputs and outcome measures as well as implementation details \citep{kerwin_making_2021}

%In a study conducted in Ethiopia, \citet{chicoine_schooling_2019} asserts that the shift to mother tongue instruction has led to a reduction in schooling and had no impact on literacy, with the negative impact being concentrated in regions that made the switch from an Amharic script to the Roman script. 

MTI has led to negative job outcomes by reducing proficiency in a dominant language \citep[e.g.,][]{angrist_effect_1997}. Having greater proficiency in a dominant language is important for success in the labor market in certain contexts \citep{chiswick_immigrant_2000,chiswick_endogeneity_2015,bleakley_language_2004,munshi_traditional_2006,lang_return_2009,aldashev_language_2009,shastry_human_2012,casale_english_2011,parinduri_effects_2018}, as it improves job search results, productivity on the job, and promotion to higher paying positions \citet{kahn_returns_2019}. %More recently, \citet{parinduri_effects_2018}, using data from Malaysia, have shown that having English as a medium of instruction improves English proficiency and educational attainment, albeit with a weaker link to later labour market outcomes. 
Hence, mother tongue instruction could negatively affect employment outcomes if it adversely affects fluency in a dominant language, although it is unclear what causes the trade-off between learning in mother tongue and mastery of a dominant language. 

%Most of the empirical analysis on the effects of MTI is conducted in contexts where an international language served as a lingua franca \citep[e.g.,][]{eriksson_does_2014,laitin_legacy_2019}. For instance, \citet{eriksson_does_2014} examined the effect of mother tongue vs. English or Afrikaans instruction using the Bantu Education Act of 1955 in South Africa as a natural experiment and finds that increasing MTI from four to six years positively affected wages, literacy, educational attainment, and English speaking skills. \citet{laitin_legacy_2019} conducted a randomized evaluation of a local language schooling program in Cameroon and found that students that were exposed to three years of the local language (Kom) in early grades scored higher in math and English tests than untreated students, who were instructed solely in English. 


%The socio-economic outcomes of shifting from a dominant language (international or domestic) to MTI are mediated through a number of factors, including the degree of people's prior exposure to the dominant language in question \citep{ramachandran_language_2017,laitin_legacy_2019}. 

The majority of existing studies on the effects of MTI are conducted in countries with an international language serving as lingua franca. Our study evaluates the impact of MTI in Ethiopia, where it was introduced as part of a larger strategy to change the character of the Ethiopian state, from a unitary system that had prioritized the assimilation of the diverse population of the country into Amharic speaking citizens, to a federation of self-governing ethnic states with their own official local languages \citep{bulcha_politics_1997,getachew_language_2006}. Arguably, the introduction of MTI was aimed at tackling the inherent contradictions of the Ethiopian state that had prioritized the development of Amharic at the expense of the other local languages. Hence, this study may yield useful insights for policy makers seeking to implement optimal policy solutions to the increasing demands of many linguistic communities around the world to learn and work in their native languages. 

%As \citet{long_bilingual_2009} rightly notes, “language as a medium of instruction confers recognition, status, and often economic benefits (e.g., teaching positions) on speakers of that language ... [and] it is ... a socio-political phenomenon that is implicated in the ongoing competition between social groups for material and symbolic resources.”

Our article is closely related to two recent studies conducted in Ethiopia \cite{seid_mother-tongue_2022} and \cite{chicoine_schooling_2019}, which, due to data constraints, are based on a less rigorous premise that the dominant ethnic language in each regional state is the only language of instruction in each state. Consequently, \cite{seid_mother-tongue_2022} excluded from its analysis all samples from the South due to the state's heterogeneity in terms of ethnicity and found that MTI has positive effects on earnings, and those effects decrease with the duration of exposure to MTI in primary school. Using a similar strategy that excludes ethnic heterogeneity within a region, \citep{chicoine_schooling_2019} examines the impact of MTI on schooling and literacy rates in Ethiopia, finding that the introduction of MTI had led to a reduction in schooling with no impact on literacy, particularly in regions where the new language of instruction was rarely used before the reform.  The data we use for our study identify the mother tongue and language of instruction of each individual, allowing us to explore the impact of variations in MTI at the individual level on their employment.


Our research has some relevance to the literature on the impact of cognitive ability on later labor market outcomes. Research indicates that there is a strong connection between educational attainment, cognitive ability, and job outcomes \citep{bowles_determinants_2001,cunha_technology_2007,heckman_fostering_2013,keuschnigg_plateauing_2023}. Cognitive skills, including problem-solving and analytical reasoning, play a crucial role in determining workplace performance \citep{deming_growing_2017,almlund_personality_2011,schmidt_general_2004}. A longitudinal study by \cite{moffitt_gradient_2011} suggests that cognitive ability in adolescence is a predictor of employment status in adulthood. Having a higher cognitive ability has also been shown to have a long term effect on career trajectories. Consistent with the findings in the literature, we demonstrate that employment opportunities vary positively in cognitive ability in labor markets that are not subject to substantial distortions. 


The paper proceeds as follows: In section \ref{section:data}, we describe the data as well as the institutional setting and the policy change that made this study possible, showing that the processes pursued by different state governments were not based on serious assessments of whether full or partial implementation of national policy would yield optimal results. After describing our key identification strategy in section \ref{section:strategy}, we proceed to discuss the main findings in section \ref{section:results}, contextualizing them with the social environment in which MTI was implemented. In section \ref{section:robustness}, we will demonstrate that our key findings remain robust by performing some tests, before concluding by highlighting the prominent policy implications of the study.
